% =====================================================================
\chapter{Teoremas límite y estimación por Monte Carlo}
\label{tema:5}
% =====================================================================

Ni la pérdida de un LLM ni su exactitud en un banco de pruebas pueden calcularse exactamente: son esperanzas respecto de una distribución desconocida y se estiman con promedios. Este tema demuestra que esos promedios convergen (ley de los grandes números), cuantifica su error (desigualdad de Chebyshev y teorema central del límite) y lo aplica al método de Monte Carlo, a los intervalos de confianza y a los minilotes del entrenamiento. Se sigue el capítulo~5 de~\cite{bt}, las lecciones~19 y~20 de~\cite{mit6041}, la clase~6b de~\cite{mit1805} y el capítulo~10 de~\cite{bh}; la §2.6 de~\cite{d2l} trata el mismo material en el contexto del aprendizaje profundo.

\section{Desigualdades de Markov y Chebyshev}

\begin{teorema}[Desigualdad de Markov]\label{teo:markov}
Si $Y\ge 0$ y $a>0$, entonces $\PP(Y\ge a)\le\E[Y]/a$.
\end{teorema}

\begin{teorema}[Desigualdad de Chebyshev]\label{teo:chebyshev}
Si $X$ tiene media $\mu$ y varianza $\sigma^{2}$, para todo $\varepsilon>0$,
\[
\PP\bigl(\abs{X-\mu}\ge\varepsilon\bigr)\le\frac{\sigma^{2}}{\varepsilon^{2}}.
\]
\end{teorema}

Las dos desigualdades valen para cualquier distribución con los momentos indicados; a cambio, suelen ser pesimistas~\cite[§5.1]{bt}.

\section{Ley débil de los grandes números}

\begin{definicion}[Convergencia en probabilidad]
Una sucesión $Y_1,Y_2,\dots$ \emph{converge en probabilidad} a la constante $c$ si, para todo $\varepsilon>0$, $\lim_{n\to\infty}\PP\bigl(\abs{Y_n-c}\ge\varepsilon\bigr)=0$.
\end{definicion}

\begin{teorema}[Ley débil de los grandes números]\label{teo:lgn}
Sean $X_1,X_2,\dots$ \iid\ con media $\mu$ y varianza finita $\sigma^{2}$. Entonces, para todo $\varepsilon>0$,
\[
\PP\bigl(\abs{\Xbar_n-\mu}\ge\varepsilon\bigr)\le\frac{\sigma^{2}}{n\varepsilon^{2}}\xrightarrow[n\to\infty]{}0,
\]
es decir, $\Xbar_n$ converge en probabilidad a $\mu$.
\end{teorema}

La ley fuerte de los grandes números afirma que, de hecho, $\Xbar_n\to\mu$ con probabilidad~1; su demostración excede este módulo~\cite[§5.5]{bt}.

\section{El método de Monte Carlo}

\begin{definicion}[Estimador de Monte Carlo]\label{def:montecarlo}
Para aproximar $I=\E[g(X)]$, se generan $X_1,\dots,X_n$ \iid\ con la distribución de $X$ y se toma $\hat I_n=\frac1n\sum_{i=1}^{n}g(X_i)$. Por el \cref{cor:media}, $\hat I_n$ es insesgado y su \emph{error típico} es $\sigma_g/\sqrt n$, donde $\sigma_g^{2}=\Var(g(X))$.
\end{definicion}

\addtocounter{definicion}{1}%

La ley $1/\sqrt n$ no depende de la dimensión del problema, lo que hace a Monte Carlo competitivo para integrales en muchas dimensiones, pero sí la hace lenta: el error se reduce a la mitad solo cuadruplicando el número de muestras.

\addtocounter{definicion}{1}%

\section{Teorema central del límite}

La ley de los grandes números dice que $\Xbar_n-\mu\to 0$; el teorema central del límite describe la forma de la distribución de ese error cuando se amplía por el factor $\sqrt n$.

\begin{teorema}[Teorema central del límite]\label{teo:tcl}
Sean $X_1,X_2,\dots$ \iid\ con media $\mu$ y varianza $\sigma^{2}\in(0,\infty)$, y sea $S_n=X_1+\dots+X_n$. Entonces, para todo $z\in\R$,
\[
\lim_{n\to\infty}\PP\Bigl(\frac{S_n-n\mu}{\sigma\sqrt n}\le z\Bigr)=\Phi(z).
\]
\end{teorema}

La demostración habitual usa funciones generatrices de momentos o funciones características~\cite[§5.4]{bt}. Lo notable es que no se exige que los $X_i$ sean normales: la normalidad es la conclusión, no la hipótesis. Así se explica la ubicuidad de la normal: el ruido térmico, los errores de medida o la suma de muchos pesos pequeños son sumas de muchos efectos independientes.

\begin{corolario}[Aproximación normal de la media]\label{cor:aprox}
Para $n$ grande, $\Xbar_n$ es aproximadamente $\N(\mu,\sigma^{2}/n)$ y, en particular,
\[
\PP\Bigl(\abs{\Xbar_n-\mu}\ge z\,\frac{\sigma}{\sqrt n}\Bigr)\approx 2\bigl(1-\Phi(z)\bigr).
\]
\end{corolario}

\begin{table}[h]
\centering\small
\caption{Tres herramientas aplicadas al \cref{ej:pi} con $n=10^{4}$ y un error de $0{,}05$, que equivale a $3{,}04$ errores típicos.}\label{tab:herramientas}
\begin{tabular}{@{}lll@{}}
\toprule
Herramienta & Afirmación & Carácter \\
\midrule
Ley de los grandes números & $\hat\pi_n\to\pi$ & cualitativa \\
Chebyshev & $\PP(\abs{\hat\pi_n-\pi}\ge 0{,}05)\le 0{,}108$ & cota rigurosa, pesimista \\
Teorema central del límite & $\PP(\abs{\hat\pi_n-\pi}\ge 0{,}05)\approx 0{,}0023$ & aproximación realista \\
\bottomrule
\end{tabular}
\end{table}

\section{Intervalos de confianza y comparación de modelos}

\begin{definicion}[Intervalo de confianza aproximado]
Sea $\hat p$ la proporción de aciertos de un modelo en $n$ preguntas independientes. El intervalo
\[
\hat p\pm 1{,}96\sqrt{\frac{\hat p(1-\hat p)}{n}}
\]
es un \emph{intervalo de confianza aproximado del $95\,\%$} para la exactitud verdadera $p$: por el \cref{cor:aprox}, contiene a $p$ con probabilidad aproximada $0{,}95$~\cite[clase~22]{mit1805}.
\end{definicion}

\addtocounter{definicion}{1}%

\section{Minilotes y gradiente estocástico}

La función de pérdida de entrenamiento es una esperanza, $L(\theta)=\E[\lambda(\theta;X)]$, donde $\lambda(\theta;x)$ es la pérdida en el ejemplo $x$ y $X$ sigue la distribución de los datos. Un \emph{minilote} es una muestra $X_1,\dots,X_B$ extraída al azar, y la pérdida del minilote es
\[
\hat L_B(\theta)=\frac1B\sum_{i=1}^{B}\lambda(\theta;X_i).
\]

\begin{proposicion}\label{prop:minilote}
Si los ejemplos del minilote son \iid\ y $\sigma^{2}=\Var(\lambda(\theta;X))$, entonces $\E[\hat L_B(\theta)]=L(\theta)$ y la desviación típica de $\hat L_B(\theta)$ es $\sigma/\sqrt B$. Si la derivada y la esperanza pueden intercambiarse, $\nabla_\theta\hat L_B$ es también un estimador insesgado de $\nabla_\theta L$, cuya covarianza decrece como $1/B$.
\end{proposicion}

\addtocounter{definicion}{1}%

El descenso de gradiente estocástico (módulo~7) avanza con estos gradientes ruidosos. Por el \cref{teo:tcl}, el ruido es aproximadamente normal, y el tamaño de lote es un compromiso entre la calidad de cada paso y su coste.

